% =========================================================
% Cardinality and Infinite Sets
% =========================================================

\section{Cardinality and Infinite Sets}

% ---------------------------------------------------------
% TOOLKIT
% ---------------------------------------------------------

\begin{exposition}
The intuitive idea of ``size'' works perfectly for finite sets: $\{a,b,c\}$
has three elements, $\{1,2,3,4\}$ has four, and any two finite sets can be
compared by counting. For infinite sets, counting breaks down, and a
different foundation is needed. The key insight, due to Cantor, is that two
sets have the ``same size'' if and only if their elements can be matched
perfectly — one-to-one and onto — without any element left over on either
side. This single idea, made precise through the notion of a bijection, gives
a theory of size that applies uniformly to both finite and infinite sets, and
produces genuinely surprising results: the integers and the rationals turn out
to have the same size as the natural numbers, while the real numbers are
provably larger.
\end{exposition}

% ---------------------------------------------------------
% Same Cardinality
% ---------------------------------------------------------

\begin{definitionbox}{Definition (Same Cardinality)}
\begin{definition}[Same Cardinality]\label{def:same-cardinality}
Two sets $A$ and $B$ have the \emph{same cardinality}, written $A \sim B$,
if there exists a bijection $f : A \to B$.
\end{definition}
\end{definitionbox}

\begin{remark*}[Standard quantified statement]
$A \sim B \;\leftrightarrow\; \exists f: A\to B,\; (\forall a_1,a_2\in A,\; f(a_1)=f(a_2)\to a_1=a_2) \wedge (\forall b\in B,\; \exists a\in A,\; f(a)=b)$.
\end{remark*}


\begin{dependencies}
\begin{itemize}
  \item \hyperref[def:bijective]{Bijective Function}
\end{itemize}
\end{dependencies}

\begin{remark*}[Same cardinality is an equivalence relation]
The relation $\sim$ is an equivalence relation on the class of all sets:
\begin{enumerate}[label=(\roman*)]
\item \emph{Reflexive}: $A \sim A$ via $\id_A$.
\item \emph{Symmetric}: if $f : A \to B$ is a bijection then $f^{-1} : B \to A$
is a bijection, so $A \sim B$ implies $B \sim A$.
\item \emph{Transitive}: if $f : A \to B$ and $g : B \to C$ are bijections then
$g \circ f : A \to C$ is a bijection, so $A \sim B$ and $B \sim C$ imply $A \sim C$.
\end{enumerate}
This is a genuine equivalence relation, so all the theory of equivalence classes
developed in the previous section applies. The equivalence class of $A$ under
$\sim$ is called the \emph{cardinality} of $A$, written $|A|$.
\end{remark*}

\begin{example*}[Infinite sets can be equinumerous with proper subsets]
\label{ex:N-even}
Let $E = \{2, 4, 6, \dots\}$ be the set of positive even integers. The
function $f : \mathbb{N} \to E$ defined by $f(n) = 2n$ is a bijection, so
$\mathbb{N} \sim E$. This seems paradoxical: $E$ is a proper subset of
$\mathbb{N}$, yet the two sets have the same cardinality. This is not a
contradiction but rather one of the defining characteristics of infinite sets
— a phenomenon that simply cannot occur for finite sets.
\end{example*}

\begin{remark*}[Dedekind's characterization of infinite sets]
The phenomenon in Example~\ref{ex:N-even} is not an accident. An infinite set
can always be matched bijectively with some proper subset of itself. This is
sometimes taken as the \emph{definition} of an infinite set (Dedekind's
definition): a set is infinite if and only if it is equinumerous with a proper
subset of itself.
\end{remark*}

% ---------------------------------------------------------
% Finite, Countable, Uncountable
% ---------------------------------------------------------

\begin{definitionbox}{Definition (Finite and Infinite Sets)}
\begin{definition}[Finite and Infinite Sets]\label{def:finite-infinite}
A set $A$ is \emph{finite} if $A = \varnothing$ or $A \sim \{1, 2, \dots, n\}$
for some $n \in \mathbb{N}$. A set is \emph{infinite} if it is not finite.
\end{definition}
\end{definitionbox}

\begin{remark*}[Standard quantified statement]
$A$ is finite $\;\leftrightarrow\; A = \varnothing \vee \exists n\in\mathbb{N},\; \exists f: A\to\{1,\dots,n\},\; f$ bijective. $A$ is infinite $\;\leftrightarrow\; \neg(A$ is finite$)$.
\end{remark*}


\begin{dependencies}
\begin{itemize}
  \item \hyperref[def:empty-set]{Empty Set}
  \item \hyperref[def:same-cardinality]{Same Cardinality}
  \item \hyperref[def:natural-numbers]{Natural Numbers}
\end{itemize}
\end{dependencies}

\begin{definitionbox}{Definition (Countable and Uncountable Sets)}
\begin{definition}[Countable and Uncountable Sets]\label{def:countable}
A set $A$ is \emph{countably infinite} if $A \sim \mathbb{N}$. A set is
\emph{countable} if it is finite or countably infinite. A set is
\emph{uncountable} if it is infinite but not countable.
\end{definition}
\end{definitionbox}

\begin{remark*}[Standard quantified statement]
$A$ is countably infinite $\;\leftrightarrow\; \exists f: A\to\mathbb{N},\; f$ bijective. $A$ is countable $\;\leftrightarrow\; A$ finite $\vee$ $A\sim\mathbb{N}$. $A$ is uncountable $\;\leftrightarrow\; A$ infinite $\wedge \neg(A\sim\mathbb{N})$.
\end{remark*}


\begin{dependencies}
\begin{itemize}
  \item \hyperref[def:finite-infinite]{Finite and Infinite Sets}
  \item \hyperref[def:same-cardinality]{Same Cardinality}
  \item \hyperref[def:natural-numbers]{Natural Numbers}
\end{itemize}
\end{dependencies}

\begin{remark*}[Countable means listable]
A set $A$ is countably infinite if and only if its elements can be arranged
into an infinite list $a_1, a_2, a_3, \dots$ with no repetitions and every
element of $A$ appearing somewhere in the list. The bijection $f : \mathbb{N}
\to A$ is precisely such a listing: $a_n = f(n)$. This is why countability is
often informally described as ``can be listed.''
\end{remark*}

% ---------------------------------------------------------
% Ordering cardinalities
% ---------------------------------------------------------

\begin{definitionbox}{Definition (Comparing Cardinalities)}
\begin{definition}[Comparing Cardinalities]\label{def:cardinality-order}
Write $|A| \leq |B|$ if there exists an injection $f : A \hookrightarrow B$.
Write $|A| < |B|$ if $|A| \leq |B|$ but $A \not\sim B$.
\end{definition}
\end{definitionbox}

\begin{remark*}[Standard quantified statement]
$|A|\leq|B| \;\leftrightarrow\; \exists f: A\to B,\; \forall a_1,a_2\in A,\; f(a_1)=f(a_2)\to a_1=a_2$ (injection exists).
\end{remark*}


\begin{dependencies}
\begin{itemize}
  \item \hyperref[def:injective]{Injective Function}
  \item \hyperref[def:same-cardinality]{Same Cardinality}
\end{itemize}
\end{dependencies}

\begin{remark*}[Why injections, not surjections]
An injection $A \hookrightarrow B$ places a ``copy'' of $A$ inside $B$
without overlap, which formalizes ``$A$ is no larger than $B$.'' The direction
matters: an injection from $A$ to $B$ is not the same as an injection from
$B$ to $A$.
\end{remark*}

% ---------------------------------------------------------
% Q countable
% ---------------------------------------------------------

\begin{theorembox}{Theorem ($\mathbb{Q}$ is countable)}
\begin{theorem}[$\mathbb{Q}$ is countable]\label{thm:Q-countable}
The set $\mathbb{Q}$ of rational numbers is countable.
\hyperref[prf:q-countable]{Go to proof.}
\end{theorem}
\end{theorembox}

\begin{dependencies}
\begin{itemize}
  \item \hyperref[def:countable]{Countable Set}
  \item \hyperref[def:natural-numbers]{Natural Numbers}
  \item \hyperref[def:bijective]{Bijective Function}
\end{itemize}
\end{dependencies}

\begin{remark*}[Proof strategy — diagonal enumeration]
The argument arranges all positive rationals into an infinite two-dimensional
array: place $p/q$ (in lowest terms, $p,q > 0$) in row $p$, column $q$. Then
traverse the array along successive diagonals — the anti-diagonal $p + q = k$
for $k = 2, 3, 4, \dots$ — skipping any fraction already seen. This produces
an explicit list of all positive rationals, and one then interleaves the
negative rationals and zero. The precise details are left as a guided proof
exercise. The core idea is: a doubly infinite list can be converted into a
singly infinite list by reading diagonals.
\end{remark*}

\begin{remark*}[What this means]
$\mathbb{Q}$ is dense in $\mathbb{R}$ — between any two real numbers there
is a rational — yet $\mathbb{Q}$ is countable. This means the rationals are
in some precise sense ``thin'' even though they are everywhere. The uncountability
of $\mathbb{R}$ that follows shows the reals are genuinely larger.
\end{remark*}

\begin{theorembox}{Theorem (Countable union of countable sets is countable)}
\begin{theorem}[Countable union of countable sets is countable]
\label{thm:countable-union}
Let $\{A_n\}_{n \in \mathbb{N}}$ be a countable family of countable sets.
Then $\bigcup_{n=1}^{\infty} A_n$ is countable.
\hyperref[prf:countable-union]{Go to proof.}
\end{theorem}
\end{theorembox}

\begin{dependencies}
\begin{itemize}
  \item \hyperref[def:countable]{Countable Set}
  \item \hyperref[def:natural-numbers]{Natural Numbers}
  \item \hyperref[def:union]{Union}
\end{itemize}
\end{dependencies}

\begin{remark*}[Proof strategy]
The argument again uses a diagonal enumeration. Arrange the elements of each
$A_n$ as a row in an infinite array (padding with repetitions if $A_n$ is
finite), then read the array along anti-diagonals. This produces a surjection
from $\mathbb{N}$ onto $\bigcup A_n$, which is enough to establish
countability. The formal proof is a standard exercise.
\end{remark*}

\begin{remark*}[Why this theorem matters downstream]
This result is used constantly. In topology: a countable union of nowhere-dense
sets is called a \emph{meager} set, and Baire's theorem says a complete metric
space is not meager — the proof requires knowing that countable unions behave
predictably. In measure theory: a countable union of sets of measure zero has
measure zero, by countable additivity. Whenever an argument says ``take the
union over all $n \in \mathbb{N}$'' and needs the result to remain manageable
in size, this theorem is the justification.
\end{remark*}

% ---------------------------------------------------------
% R uncountable
% ---------------------------------------------------------

\begin{theorembox}{Theorem ($\mathbb{R}$ is uncountable)}
\begin{theorem}[$\mathbb{R}$ is uncountable]\label{thm:R-uncountable}
The set $\mathbb{R}$ of real numbers is not countable.
\hyperref[prf:r-uncountable]{Go to proof.}
\end{theorem}
\end{theorembox}

\begin{dependencies}
\begin{itemize}
  \item \hyperref[def:countable]{Countable Set}
  \item \hyperref[def:finite-infinite]{Finite and Infinite Sets}
  \item \hyperref[thm:countable-union]{Countable Union of Countable Sets}
\end{itemize}
\end{dependencies}

\begin{remark*}[Proof strategy — Cantor's diagonal argument]
This is one of the most important proof techniques in all of mathematics.
Suppose, for contradiction, that $\mathbb{R}$ were countable, so its elements
could be listed: $r_1, r_2, r_3, \dots$. Write each $r_n$ in decimal
expansion. Construct a new real number $x$ by choosing the $n$-th decimal
digit of $x$ to differ from the $n$-th decimal digit of $r_n$ (and avoiding
$9$ to prevent the $0.999\dots = 1$ ambiguity). Then $x \neq r_n$ for every
$n$, because $x$ and $r_n$ differ in the $n$-th decimal place. So $x$ is not
on the list — contradicting the assumption that the list contains every real
number.

The key move is the diagonal construction: to defeat every row $r_n$
simultaneously, make $x$ disagree with row $n$ at position $n$. This
``diagonal sweep'' is a technique that reappears in logic (Gödel's
incompleteness theorem), computability theory (the halting problem), and the
proof of Cantor's theorem below.
\end{remark*}

\begin{remark*}[Consequence]
Since $\mathbb{Q}$ is countable and $\mathbb{R}$ is uncountable, the
irrational numbers $\mathbb{R} \setminus \mathbb{Q}$ must be uncountable:
if they were countable, then $\mathbb{R} = \mathbb{Q} \cup (\mathbb{R}
\setminus \mathbb{Q})$ would be a countable union of countable sets, hence
countable — contradiction.
\end{remark*}

% ---------------------------------------------------------
% Cantor's theorem
% ---------------------------------------------------------

\begin{theorembox}{Theorem (Cantor's Theorem)}
\begin{theorem}[Cantor's Theorem]\label{thm:cantor}
For any set $A$, there is no surjection $A \to \mathcal{P}(A)$. In
particular, $|A| < |\mathcal{P}(A)|$.
\hyperref[prf:cantor]{Go to proof.}
\end{theorem}
\end{theorembox}

\begin{remark*}[Standard quantified statement]
$\forall A:\; \neg\exists f: A\to\mathcal{P}(A),\; \forall S\in\mathcal{P}(A),\; \exists a\in A,\; f(a)=S$. Equivalently: $\forall f: A\to\mathcal{P}(A),\; \exists D\in\mathcal{P}(A),\; \forall a\in A,\; f(a)\neq D$.
\end{remark*}


\begin{dependencies}
\begin{itemize}
  \item \hyperref[def:function]{Function}
  \item \hyperref[def:surjective]{Surjective Function}
  \item \hyperref[def:power-set]{Power Set}
  \item \hyperref[ax:separation]{Axiom Schema of Separation}
\end{itemize}
\end{dependencies}

\begin{remark*}[Proof strategy — the same diagonal idea]
Suppose $f : A \to \mathcal{P}(A)$ is any function. Construct the set
\[
D \;=\; \{\, a \in A \mid a \notin f(a) \,\}.
\]
$D$ is a well-defined subset of $A$, so $D \in \mathcal{P}(A)$. Claim: $D$
is not in the image of $f$. For suppose $D = f(a_0)$ for some $a_0 \in A$.
Then:
\begin{itemize}
\item if $a_0 \in D$, then by definition $a_0 \notin f(a_0) = D$ — contradiction.
\item if $a_0 \notin D$, then by definition $a_0 \in f(a_0) = D$ — contradiction.
\end{itemize}
Either case is impossible, so $D \neq f(a_0)$ for every $a_0 \in A$, and $f$
is not surjective. Since $f$ was arbitrary, no surjection $A \to
\mathcal{P}(A)$ exists.

Notice the structure: $D$ is defined by ``$a \in D$ iff $a \notin f(a)$,''
which is a diagonal construction — the membership of $a$ in $D$ is decided
by looking at position $(a,a)$ in the ``membership table'' of $f$, and then
disagreeing. This is the abstract form of the same move used in the proof
that $\mathbb{R}$ is uncountable.
\end{remark*}

\begin{remark*}[A hierarchy of infinities]
Cantor's theorem implies there is no largest cardinality: for any set $A$,
the power set $\mathcal{P}(A)$ is strictly larger. Starting from $\mathbb{N}$,
the sets
\[
\mathbb{N},\quad
\mathcal{P}(\mathbb{N}),\quad
\mathcal{P}(\mathcal{P}(\mathbb{N})),\quad \dots
\]
form a strictly increasing sequence of cardinalities. One can show
$|\mathcal{P}(\mathbb{N})| = |\mathbb{R}|$ via the characteristic function
bijection (described below), so the second level is already the cardinality
of the continuum.
\end{remark*}

% ---------------------------------------------------------
% Characteristic functions and B^A
% ---------------------------------------------------------

\begin{definitionbox}{Definition (Function Space $B^A$)}
\begin{definition}[Function Space $B^A$]\label{def:function-space}
For sets $A$ and $B$, the \emph{function space} $B^A$ denotes the set of
all functions from $A$ to $B$:
\[
B^A \;:=\; \{\, f : A \to B \,\}.
\]
\end{definition}
\end{definitionbox}

\begin{dependencies}
\begin{itemize}
  \item \hyperref[def:function]{Function}
  \item \hyperref[def:set-membership]{Set and Membership}
\end{itemize}
\end{dependencies}

\begin{remark*}[Why the exponential notation]
For finite sets with $|A| = m$ and $|B| = n$, the number of functions
$A \to B$ is $n^m$ (there are $n$ choices for each of the $m$ inputs).
The notation $B^A$ reflects this count. In particular, $|\{0,1\}^A| = 2^{|A|}$
for finite $A$, which matches $|\mathcal{P}(A)| = 2^{|A|}$.
\end{remark*}

\begin{remark*}[The characteristic function bijection]
\label{rem:characteristic-fn}
For any set $A$, the power set $\mathcal{P}(A)$ and the function space
$\{0,1\}^A$ have the same cardinality via the \emph{characteristic function}
bijection:
\[
\mathcal{P}(A) \;\to\; \{0,1\}^A,
\qquad
S \;\mapsto\; \mathbf{1}_S,
\]
where $\mathbf{1}_S(a) = 1$ if $a \in S$ and $\mathbf{1}_S(a) = 0$ if
$a \notin S$. This bijection is why Cantor's theorem is sometimes stated as:
there is no surjection $A \to \{0,1\}^A$. The notation $2^A$ is used
interchangeably with $\{0,1\}^A$, and $|\mathcal{P}(A)| = 2^{|A|}$ is
the general statement.

In functional analysis, spaces of the form $B^A$ — functions from $A$ to
$B$ — are the central objects of study. Knowing they are sets (with a definite
cardinality when $A$ and $B$ are concrete) is the set-theoretic foundation
for everything that follows.
\end{remark*}

% ---------------------------------------------------------
% Schröder–Bernstein
% ---------------------------------------------------------

\begin{theorembox}{Theorem (Schröder--Bernstein Theorem)}
\begin{theorem}[Schröder--Bernstein Theorem]\label{thm:schroder-bernstein}
If $|A| \leq |B|$ and $|B| \leq |A|$, then $A \sim B$.

Equivalently: if there exist injections $f : A \hookrightarrow B$ and
$g : B \hookrightarrow A$, then there exists a bijection $h : A \to B$.
\hyperref[prf:schroder-bernstein]{Go to proof.}
\end{theorem}
\end{theorembox}

\begin{remark*}[Standard quantified statement]
$\bigl(\exists f: A\hookrightarrow B\bigr) \wedge \bigl(\exists g: B\hookrightarrow A\bigr) \;\to\; \exists h: A\to B,\; h$ bijective.
\end{remark*}


\begin{dependencies}
\begin{itemize}
  \item \hyperref[def:cardinality-order]{Comparing Cardinalities}
  \item \hyperref[def:injective]{Injective Function}
  \item \hyperref[def:bijective]{Bijective Function}
\end{itemize}
\end{dependencies}

\begin{remark*}[Why this theorem is indispensable]
In practice, proving $A \sim B$ by exhibiting an explicit bijection can be
very difficult. Schröder--Bernstein gives a two-injection strategy: find any
injection each way, and the theorem guarantees a bijection exists (even if
you cannot describe it explicitly). This is a powerful tool for cardinality
computations. For example, to show $|(0,1)| = |\mathbb{R}|$, find an
injection $(0,1) \hookrightarrow \mathbb{R}$ (inclusion) and an injection
$\mathbb{R} \hookrightarrow (0,1)$ (e.g., $x \mapsto \frac{1}{\pi}
\arctan(x) + \frac{1}{2}$); the theorem does the rest.
\end{remark*}

\begin{remark*}[Proof idea]
The proof constructs $h$ by partitioning $A$ into two pieces: those elements
whose ``ancestry'' under repeated application of $g \circ f$ is finite (they
go through $f$), and those whose ancestry is infinite (they go through
$g^{-1}$). The formal argument is a standard exercise; the key insight is that
the two injections together cover $A$ in a structured enough way to stitch
together a bijection.
\end{remark*}

\begin{remark*}[Sequences as functions]
\label{rem:sequences-as-functions}
A \emph{sequence} in a set $X$ is a function $f : \mathbb{N} \to X$. The
element $f(n)$ is written $x_n$, and the sequence is written $(x_n)_{n \in
\mathbb{N}}$ or simply $(x_n)$. This is not a new object: it is an element
of the function space $X^{\mathbb{N}}$. When metric spaces and topology
discuss convergence of sequences $x_1, x_2, x_3, \dots$ in an arbitrary space
$X$, the formal object is always a function $\mathbb{N} \to X$, and all the
function machinery — injectivity, composition, preimages — applies to it
directly.
\end{remark*}

\begin{remark*}[Transition]
Cardinality distinguishes between the countable (manageable, listable) and the
uncountable (genuinely larger). This distinction organizes much of what
follows. In metric spaces: separable spaces have countable dense subsets, and
separability is a strong structural property. In measure theory: $\sigma$-algebras
are closed under \emph{countable} unions — not arbitrary ones — precisely because
countable operations preserve the kind of control that makes measure theory work.
In functional analysis: $L^2$ spaces have countable orthonormal bases
(Hilbert spaces are separable), while non-separable spaces behave very
differently. The line between countable and uncountable is not a technicality;
it is one of the organizing principles of analysis.
\end{remark*}
